\chapter{From Newton to Hamilton}\label{ch:classical}

\begin{etymology}
\textbf{Mechanics} comes from the Greek \emph{mēkhanikē}, the art of machines,
from \emph{mēkhanē}, ``machine'' or ``contrivance''. It came to mean the study of
how things move under fixed rules.

\textbf{Energy} comes from the Greek \emph{enérgeia}: \emph{en-}, ``in'', and
\emph{érgon}, ``work''. Something at work, in action.

\textbf{Hamiltonian} is named after William Rowan Hamilton (1805--1865), the Irish
mathematician who rewrote Newton's mechanics in terms of position and momentum.
\end{etymology}

Quantum mechanics was not invented from scratch. It grew out of classical
mechanics at the moment classical mechanics stopped matching experiments, and it
kept a surprising amount of the old structure. To see what changed, you need to
know what was there before.

Classical mechanics was written three times, by Newton, Lagrange and Hamilton.
The three versions say the same thing about the physical world. What changes is
the language, and each new language is more abstract and more powerful than the
last. Hamilton's is the one quantum mechanics kept.

\section{Newton: forces and trajectories}

\begin{history}{1687 \textperiodcentered{} Cambridge}
Isaac Newton published the \emph{Philosophiæ Naturalis Principia Mathematica} in
1687, after Edmond Halley pushed him to write up his work on orbits and paid for
the printing. The three laws of motion and the law of gravitation explained, with
one set of equations, why apples fall and why planets follow ellipses.
\end{history}

The second law, $F = ma$, says that force changes motion. Written out, the
acceleration is the rate of change of the rate of change of the position,
\[
  m\,\frac{d^2x}{dt^2} = F .
\]
Notice what this equation does not say. It does not give the position. It gives a
relation between the position and how fast it changes, and the position itself has to
be found from that relation. An equation of this kind is a \emph{differential
equation}, and Newton had to invent the mathematics to write it.

\begin{history}{1665--1736 \textperiodcentered{} Woolsthorpe, Hanover and London}
In 1665 and 1666, with Cambridge closed by the plague, the young Newton worked out at
his family's farm in Woolsthorpe what he called the \emph{method of fluxions}: a way
of computing rates of change (fluxions) of quantities that flow in time. In a treatise
written around 1671 he sorted the equations that relate a quantity to its fluxions into
types and showed how to solve some of them, but he kept it unpublished; it appeared
only in 1736, after his death. Gottfried Wilhelm Leibniz, in Hanover, found the same
calculus independently, published first in 1684, and gave it the notation $dx/dt$ that
is still used. He also coined the name, \emph{aequatio differentialis}. The quarrel over
who came first poisoned English and Continental mathematics for decades.
\end{history}

\begin{concept}{Differential equation}
An equation that relates an unknown function to its rates of change. The equation
only says how the system changes from one instant to the next. Solving it means
finding the function itself: the one that says where the system will be at every later
moment, how it evolves. That function is the prize, and finding it is the whole point of
writing the equation. When the function depends on a single variable, usually time, it
is an \emph{ordinary} differential equation (ODE): the spring below, and every equation
of motion in this chapter. When it depends on several variables at once, such as
position and time, the rates of change are \emph{partial} derivatives and it is a
\emph{partial} differential equation (PDE): the vibrating string (d'Alembert, 1747), the
heat equation (Fourier, 1807, Chapter~\ref{ch:fourier}) and Schrödinger's equation
(Chapter~\ref{ch:quantum}).

\textbf{What it fixes:} physical laws are local. They say what happens in the next
instant, here, given what is happening now. A differential equation is exactly that
kind of statement, and it predicts the whole future from it.

\textbf{What it costs:} most differential equations have no formula for their
solution. You either find one of the few that do, or you compute the solution step by
step.
\end{concept}

The state of a particle is its position and velocity, $(x, v)$. Give the state at one
instant and the forces, and the differential equation fixes the whole future. This is
classical determinism.

The standard example is a mass on a spring. The spring pulls back with a force
proportional to the displacement, $F = -kx$, so the equation of motion is the ODE
$m\ddot{x} = -kx$ (a dot means a derivative in time, Newton's own notation). Solving it
means finding a function whose second derivative is $-k/m$ times itself, and the
cosine is such a function:
\[
  m\ddot{x} = -kx, \qquad x(t) = A\cos(\omega t + \phi), \qquad
  \omega = \sqrt{k/m}.
\]
The mass oscillates forever with frequency $\omega$. This \emph{harmonic
oscillator} comes back again and again in physics, and it comes back in
quantum computing too: a superconducting qubit is a harmonic oscillator with a
small defect added on purpose (Chapter~\ref{ch:hardware}).

\begin{figure}[htbp]
\centering
\begin{tikzpicture}[x=1mm, y=1mm]
  % equilibrium line
  \draw[black!35, dashed] (41,8) -- (41,-39) node[below, hf/note] {rest position, $x = 0$};
  \foreach \row/\d/\lab in {0/12/a, 1/0/b, 2/-12/c} {
    \pgfmathsetmacro{\y}{-\row*16}
    \pgfmathsetmacro{\bx}{36+\d}
    \pgfmathsetmacro{\seg}{(\bx-6)/9}
    % wall and floor
    \draw[line width=1pt] (0,\y-5) -- (0,\y+5);
    \foreach \h in {-4,-2,...,4} \draw[black!45] (0,\y+\h) -- (-2,\y+\h-2);
    \draw[black!40] (0,\y-5) -- (64,\y-5);
    % spring
    \draw[line width=0.7pt] (0,\y) -- (2,\y);
    \draw[line width=0.7pt, decorate, decoration={coil, aspect=0.5,
      segment length=\seg mm, amplitude=2.2mm}] (2,\y) -- (\bx-3,\y);
    \draw[line width=0.7pt] (\bx-3,\y) -- (\bx,\y);
    % block
    \filldraw[fill=hfConcept!15, draw=hfConcept, line width=0.8pt]
      (\bx,\y-5) rectangle (\bx+10,\y+5);
    \node[font=\small] at (\bx+5,\y) {$m$};
    \node[font=\sffamily\small\bfseries, text=black!60] at (-6,\y) {(\lab)};
  }
  % forces
  \draw[hfAlert, line width=1pt, -{Stealth[length=2.2mm]}] (58,7) -- (47,7)
    node[midway, above, hf/note, text=hfAlert] {$F = -kA$};
  \node[hf/note, anchor=west] at (62,0) {$x = +A$};
  \node[hf/note, anchor=west] at (62,-16) {$x = 0$};
  \draw[hfAlert, line width=1pt, -{Stealth[length=2.2mm]}] (24,-25) -- (35,-25)
    node[midway, above, hf/note, text=hfAlert] {$F = +kA$};
  \node[hf/note, anchor=west] at (62,-32) {$x = -A$};
  % x(t)
  \begin{axis}[at={(82mm,-37mm)}, anchor=south west, width=52mm, height=50mm,
      xmin=0, xmax=1.6, ymin=-1.3, ymax=1.3,
      xtick={0,0.25,0.5,1}, xticklabels={$0$,$\frac{T}{4}$,$\frac{T}{2}$,$T$},
      ytick={-1,0,1}, yticklabels={$-A$,$0$,$+A$},
      xlabel={$t$}, ylabel={$x(t)$},
      tick label style={font=\scriptsize}, label style={font=\small},
      axis lines=middle, axis line style={black!50},
      xlabel style={at={(ticklabel* cs:1.02)}, anchor=west},
      ylabel style={at={(ticklabel* cs:1.02)}, anchor=south}]
    \addplot[hfConcept, line width=0.9pt, domain=0:1.55, samples=120] {cos(360*x)};
    \addplot[only marks, mark=*, mark size=1.6pt, hfAlert] coordinates {(0,1) (0.25,0) (0.5,-1)};
    \node[font=\sffamily\scriptsize\bfseries, text=black!60, anchor=south west] at (axis cs:0.02,1) {(a)};
    \node[font=\sffamily\scriptsize\bfseries, text=black!60, anchor=south east] at (axis cs:0.23,0.02) {(b)};
    \node[font=\sffamily\scriptsize\bfseries, text=black!60, anchor=north west] at (axis cs:0.52,-1) {(c)};
  \end{axis}
\end{tikzpicture}
\caption{A mass on a spring. (a) Stretched by $A$, the mass stops for an instant and
is pulled back with force $-kA$. (b) At the rest position the force is zero, but the
mass is moving at its fastest and overshoots. (c) Compressed by $A$, it stops again and
is pushed out. The position
traces $x(t) = A\cos\omega t$, and one full cycle takes $T = 2\pi/\omega$.}\label{fig:spring}
\end{figure}

Newton's method works well when you can write down every force. It becomes painful
when the system has constraints: a pendulum on a rod, a bead on a wire, a ball rolling
on a surface. What hurts is writing the equations of motion, and the pendulum shows
why.

A mass $m$ hangs from a rigid rod of length $\ell$ and swings. Newton's law works with
forces and with the coordinates $x$ and $y$ of the mass, so you write $F = ma$ once for
each direction. Two forces act: gravity, $mg$, which you know, and the tension $T$ of
the rod, which you do not. The tension is whatever it takes to keep the mass on its
circle, and it changes all the time as the mass swings. With the pivot at the origin,
\[
  m\ddot x = -T\,\frac{x}{\ell}, \qquad
  m\ddot y = -T\,\frac{y}{\ell} - mg, \qquad
  x^2 + y^2 = \ell^2 .
\]
Three unknown functions, $x(t)$, $y(t)$ and $T(t)$, two differential equations and a
constraint. You have to use the constraint to eliminate the tension by hand before you
can solve anything, and the tension is a quantity you never wanted to know. For a
single pendulum this is a page of algebra. For a double pendulum, or a machine with a
dozen joints, it becomes unmanageable.

\begin{figure}[htbp]
\centering
\begin{tikzpicture}[x=1mm, y=1mm]
  \fill[black!15] (-12,0) rectangle (12,2.5);
  \draw[black!60, line width=0.8pt] (-12,0) -- (12,0);
  \draw[black!30, dashed] (0,0) -- (0,-36);
  \draw[black!40] (0,-36) arc (-90:-55:36);
  \draw[black!40] (0,-36) arc (-90:-125:36);
  % rod at theta = 28 degrees
  \coordinate (m) at ({36*sin(28)},{-36*cos(28)});
  \draw[black!70, line width=1.2pt] (0,0) -- (m);
  \fill[black!70] (0,0) circle (0.9);
  \draw[black!60] (0,-10) arc (-90:-62:10);
  \node[font=\small] at (2.6,-12.5) {$\theta$};
  \node[font=\small, left] at ($(0,0)!0.55!(m)$) {$\ell$};
  % tension and weight
  \draw[hfAlert, line width=1pt, -{Stealth[length=2.2mm]}] (m) -- ($(m)!11mm!(0,0)$)
    node[pos=0.8, left=2pt, hf/note, text=hfAlert] {$T$ (unknown)};
  \draw[hfConcept, line width=1pt, -{Stealth[length=2.2mm]}] (m) -- ++(0,-12)
    node[below, hf/note, text=hfConcept] {$mg$};
  \filldraw[fill=hfConcept!15, draw=hfConcept, line width=0.8pt] (m) circle (2.6);
  \node[font=\small] at (m) {$m$};
  % axes
  \draw[black!45, -{Stealth[length=1.8mm]}] (22,-30) -- (30,-30) node[right, hf/note] {$x$};
  \draw[black!45, -{Stealth[length=1.8mm]}] (22,-30) -- (22,-22) node[above, hf/note] {$y$};
  % comparison
  \node[hf/note, anchor=west, align=left, text width=58mm] at (38,-8)
    {\textbf{Newton:} unknowns $x$, $y$ and $T$; two equations plus the constraint
     $x^2 + y^2 = \ell^2$; eliminate $T$ by hand.};
  \node[hf/note, anchor=west, align=left, text width=58mm] at (38,-26)
    {\textbf{Lagrange:} one coordinate, $\theta$; one equation,
     $\ddot\theta = -(g/\ell)\sin\theta$; $T$ never appears.};
\end{tikzpicture}
\caption{A pendulum. Newton's law needs every force, including the tension of the
rod, which is unknown. Lagrange's method uses the angle, which the rod allows to change
freely, and the tension drops out.}\label{fig:pendulum}
\end{figure}

\section{Lagrange: energy instead of forces}

\begin{history}{1788 \textperiodcentered{} Paris}
Joseph-Louis Lagrange published the \emph{Mécanique analytique} in 1788. Its
preface boasts that the book contains no figures: mechanics had become pure
analysis, operations on equations. The book starts from a principle, not from
forces: the principle of virtual work, which Jean d'Alembert had extended to bodies in
motion in 1743. The principle of least action, which Pierre-Louis Maupertuis had
stated vaguely in 1744 and Leonhard Euler and Lagrange had made precise, comes out of
it in the book as a consequence.
\end{history}

Newton and Lagrange looked at the same pendulum and saw different things. Newton saw a
body driven by forces: list every force, add them up, and the sum gives the
acceleration. Lagrange saw that, under those same forces, the motion obeys a
principle: of all the ways the system could go from here to there, it takes the one
that satisfies a single condition. Newton described the system with the forces, and
Lagrange described it with the principle. The forces do not vanish in Lagrange's
version; they are folded into the energies. Define the \emph{Lagrangian} as kinetic
energy minus potential energy, $L = T - V$. The path the system actually takes is
the one that makes the \emph{action} $S = \int L\,dt$ stationary.

\begin{concept}{Principle of least action}
Fix where the system starts and where it ends, and when. Between those two moments it
could, in imagination, take any path at all. For each imaginable path, add up $T - V$
along the way; that sum is the action,
\[
  S = \int_{t_1}^{t_2} (T - V)\,dt .
\]
The path nature actually takes is the one where $S$ is \emph{stationary}: nudge the
path slightly in any direction and, to first order, $S$ does not change. Usually that
is a minimum, which is where the word \emph{least} comes from, but strictly it is only
stationary.

\begin{center}
\begin{tikzpicture}[x=1mm, y=1mm]
  \draw[black!45, -{Stealth[length=1.8mm]}] (0,0) -- (62,0) node[right, hf/note] {$t$};
  \draw[black!45, -{Stealth[length=1.8mm]}] (0,0) -- (0,30) node[above, hf/note] {$q$};
  \coordinate (A) at (6,6);
  \coordinate (B) at (54,20);
  \draw[black!35, dashed, line width=0.6pt] (A) .. controls (20,30) and (35,-4) .. (B);
  \draw[black!35, dashed, line width=0.6pt] (A) .. controls (18,-6) and (40,34) .. (B);
  \draw[black!35, dashed, line width=0.6pt] (A) .. controls (30,8) and (34,6) .. (B);
  \draw[hfConcept, line width=1.2pt] (A) .. controls (22,22) and (38,26) .. (B);
  \fill[black!70] (A) circle (1) node[above left, hf/note] {start, $t_1$};
  \fill[black!70] (B) circle (1) node[below right, hf/note] {end, $t_2$};
  \node[hf/note, text=hfConcept, anchor=south] at (30,25) {the path taken: $S$ stationary};
  \node[hf/note, anchor=west] at (64,9) {other imaginable paths};
  \draw[black!35] (63.5,9) -- (45,9.5);
\end{tikzpicture}
\end{center}

Pierre-Louis Maupertuis proposed in 1744 that nature always spends the least
``action'', a vague quantity he partly justified by theology. Euler made it precise the
same year for a single particle, and Lagrange, in 1760, extended it to any system.

\textbf{What it fixes:} instead of following forces instant by instant, you compare whole
histories using one number. That number is the same in any coordinates, so you can
choose the convenient ones, and forces that do no work never appear.

\textbf{What it costs:} it is stated with both ends fixed, as if the system knew where it
was going. It does not; the Euler--Lagrange equations turn it back into a rule for the
next instant. And it only applies, in this simple form, to systems that do not lose
energy to friction.
\end{concept}

From that one principle come the Euler--Lagrange equations:
\[
  \frac{d}{dt}\frac{\partial L}{\partial \dot{q}} - \frac{\partial L}{\partial q} = 0 .
\]
Read it as a balance. $\partial L/\partial\dot q$ is how much $L$ depends on the
velocity; for the usual kinetic energy it is the momentum, and the first term is how
fast that momentum changes. $\partial L/\partial q$ is how much $L$ depends on the
position; for the usual potential energy it is the force. The equation says that at
every instant the two are equal: the momentum changes exactly as fast as the force
pushes. That is Newton's law again, but $q$ can now be any quantity you like, a
distance or an angle, and choosing it well is the whole trick. It is an ordinary
differential equation, one for each coordinate, because the only variable is time.

Nothing in it says the system has to be mechanical. All it needs is a Lagrangian: a
function of some quantity and of its rate of change, whose integral the system makes
stationary. In an electric circuit the quantity is the charge on the capacitor. The
energy in the coil plays the part of kinetic energy, the energy in the capacitor
plays the part of potential energy, and the same equation gives the circuit's
equation; Chapter~\ref{ch:hardware} starts from the energy of such a circuit to
describe a superconducting qubit. For a field, such as the electromagnetic field,
there is one coordinate at every point of space, the equation becomes a partial
differential equation, and out come Maxwell's equations. David Hilbert obtained the
equations of general relativity this way in 1915, and the Standard Model of particle
physics is written today as a single Lagrangian. What does not fit, in this simple
form, is friction: a system that loses energy has no action to make stationary.

For the pendulum of Figure~\ref{fig:pendulum}, Lagrange chooses the angle $\theta$
instead of $x$ and $y$, and two things disappear. The rod's constraint disappears,
because any angle keeps the bob at distance $\ell$. The tension disappears, because
it pulls along the rod while the bob moves across it, so it does no work and never
enters the energies. What is left is a single equation:
\[
  \ddot\theta = -\frac{g}{\ell}\sin\theta .
\]
For small swings it is the spring's equation, and its solution is again a cosine,
$\theta(t) = \theta_0\cos\bigl(\sqrt{g/\ell}\,t\bigr)$: the function that says where
the pendulum is at every moment.

\section{Hamilton: position and momentum}

\begin{history}{1827--1835 \textperiodcentered{} Dublin}
William Rowan Hamilton was made Astronomer Royal of Ireland in 1827, at 21, while
still an undergraduate. His first big work was not about mechanics at all but about
light. In 1827 he read his \emph{Theory of Systems of Rays} to the Royal Irish
Academy: a whole family of light rays, passing through lenses and mirrors, could be
described by a single function. Optics already had its own principle of this kind,
Pierre de Fermat's from 1662: light goes from one point to another along the path
that takes the least time. Hamilton noticed that mechanics, with Maupertuis's least
action, had the same shape, and in two papers of 1834 and 1835 he carried his optical
method over to moving bodies.
\end{history}

So Hamilton was not fixing a failure of Lagrange. Lagrange's equations solved every
problem of mechanics that Hamilton's would solve, and in the nineteenth century
nobody computed an orbit more easily because of him. What Hamilton was after was the
analogy: one theory for rays of light and for the paths of bodies. Following it, he
ended up with a different way of writing the state of a system, and that turned out
to matter more than the analogy itself.

In Lagrange, the unknown is the position $q(t)$, and the velocity is only its
derivative. To know where the system goes next you need both, so you solve one
second-order equation. Hamilton promoted the momentum to an unknown of its own,
with the same standing as the position. The price is twice as many equations; the
gain is that each one is first-order, and that the pair $(q, p)$ at one instant is
the whole state of the system, nothing missing.

\begin{concept}{Momentum}
Momentum is mass times velocity, $p = mv$, and it measures how much motion a body
carries: a truck and a bicycle at the same speed do not stop the same way. René
Descartes called it the quantity of motion in 1644, but he multiplied size by speed
and ignored direction, and his rules for collisions came out wrong. In 1668 and 1669,
answering a call from the Royal Society, John Wallis, Christopher Wren and Christiaan
Huygens showed that direction counts: momenta add as arrows, and in a collision their
total does not change. Newton put it at the center of his second law in 1687: force is
the rate at which momentum changes.

Since then the idea has stretched far beyond mass times velocity. In Lagrange and
Hamilton the momentum of a coordinate is $p = \partial L/\partial\dot q$, whatever the
coordinate is; for an angle it is the angular momentum. Light has no mass and still
carries momentum, $p = E/c$: James Clerk Maxwell predicted in 1873 that light pushes on
what it hits, and Pyotr Lebedev measured the push in 1900. In 1924 Louis de Broglie
went the other way and gave every particle a wave, with a wavelength set by its
momentum, $\lambda = h/p$. In 1918 Emmy Noether showed what all of these have in
common: momentum is the quantity that is conserved because the laws of physics are
the same here as one meter to the left. In quantum mechanics it becomes the operator
that moves a state through space.

\textbf{What it fixes:} it gives one conserved quantity to follow through any
collision, and in Hamilton's mechanics it completes the state: position and momentum
together fix the whole future.

\textbf{What it costs:} the picture of mass times velocity stops working quickly. The
momentum of an angle has different units, the momentum of light has no mass in it,
and in quantum mechanics position and momentum can no longer both be sharp.
\end{concept}

Write $q$ for the coordinate, whatever it is, and $p$ for its momentum,
$p = \partial L/\partial\dot q$. For a mass on a spring $p = m\dot x$; for a pendulum
it is the angular momentum $m\ell^2\dot\theta$. Hamilton then built a new function,
the \emph{Hamiltonian}:
\[
  H(q, p) = p\dot{q} - L .
\]
Read it this way: the Hamiltonian of $q$ and $p$ is the total energy of the system,
written as a function of where the system is and how much momentum it carries. The
Lagrangian is $T - V$, written with position and velocity; the Hamiltonian is
$T + V$, written with position and momentum. The term $p\dot q$ is only the
bookkeeping that makes the swap. For the spring, $p\dot x = m\dot x^2 = 2T$, so
$H = 2T - (T - V) = T + V$. Swapping a variable for a slope in this way is older than
Hamilton: Adrien-Marie Legendre used it in 1787, and it is called the Legendre
transform.

The motion becomes two first-order equations, symmetric in $q$ and $p$:
\[
  \dot q = \frac{\partial H}{\partial p}, \qquad
  \dot p = -\frac{\partial H}{\partial q}.
\]
The first says the position changes at the rate set by how the energy depends on the
momentum; for the spring, $\partial H/\partial p = p/m$, which is the velocity. The
second says the momentum changes at the rate set by how the energy depends on the
position, with a minus sign; for the spring, $-\partial H/\partial x = -kx$, which is
the force. The minus sign is what pushes the system downhill in potential energy.

The state of the system is a point $(q, p)$ in \emph{phase space}, and time moves
that point along a flow. For the spring, $H = p^2/2m + kx^2/2$. Energy is
conserved, so each trajectory stays on a curve of constant $H$, and for this $H$ the
curve is an ellipse.

\begin{figure}[htbp]
\centering
\begin{tikzpicture}[scale=0.9]
  \draw[black!40, -{Stealth[length=2mm]}] (-3.4,0) -- (3.6,0) node[right, font=\sffamily\small, text=black] {$x$};
  \draw[black!40, -{Stealth[length=2mm]}] (0,-2.6) -- (0,2.8) node[above, font=\sffamily\small, text=black] {$p$};
  \foreach \r/\c in {0.7/hfConcept, 1.2/hfConcept, 1.7/hfBus, 2.2/hfHistory} {
    \draw[\c, line width=0.9pt] (0,0) ellipse ({1.35*\r} and {\r});
  }
  \draw[hfHistory, line width=0.9pt, -{Stealth[length=2.4mm]}] (0.05,2.2) -- (0.3,2.18);
  \draw[black!35] (3.05,1.75) -- (2.3,1.45);
  \node[hf/note, anchor=west] at (3.1,1.75) {higher energy};
  \draw[black!35] (3.05,-1.75) -- (0.72,-0.5);
  \node[hf/note, anchor=west] at (3.1,-1.75) {lower energy};
  \fill[hfHistory] (2.97,0) circle (2.2pt) node[above right, hf/note, text=hfHistory] {(a)};
  \fill[hfHistory] (0,2.2) circle (2.2pt) node[above left, hf/note, text=hfHistory] {(b)};
  \fill[hfHistory] (-2.97,0) circle (2.2pt) node[above left, hf/note, text=hfHistory] {(c)};
\end{tikzpicture}
\caption{Phase space of the harmonic oscillator. Each ellipse is one trajectory,
at one fixed energy, and the system goes around its ellipse forever, clockwise. The
points (a), (b) and (c) are the three moments of Figure~\ref{fig:spring}: stretched
and at rest, passing the rest position at full speed, compressed and at
rest.}\label{fig:phase}
\end{figure}

\begin{reflection}{Is a point of this diagram an energy, or a rate of change?}
Neither. A point of the diagram is a state: one position and one momentum at the
same instant. The point (a) is the spring stretched with the mass at rest; the point
(b) is the mass passing the rest position at full speed.

The energy is not on either axis. It is a number attached to each point, $H(x, p)$,
the way altitude is attached to each point of a map, and the ellipses are the
contour lines of that map. Every point on one ellipse has the same energy, and points
on bigger ellipses have more. At (a) all of it is potential energy, because $p = 0$;
at (b) all of it is kinetic, because $x = 0$; in between it is shared, and the sum
stays the same.

The rate of change is a third thing: how fast, and in which direction, the point
moves. Hamilton's equations give it, $\dot x = \partial H/\partial p$ and
$\dot p = -\partial H/\partial x$, and those are the slopes of the energy map turned
by a right angle. Someone who always walks at right angles to the slope of a hill
follows a contour line and never climbs or descends; that is why the system stays on
its ellipse. So the energy is not the rate of change. The energy is the height of the
map, and the rate of change comes from how steep the map is.
\end{reflection}

Hamilton's form has one more piece that matters later. For any two quantities
$f$ and $g$ of the system, define the \emph{Poisson bracket}, introduced by
Siméon Denis Poisson in 1809:
\[
  \{f, g\} = \frac{\partial f}{\partial q}\frac{\partial g}{\partial p}
           - \frac{\partial f}{\partial p}\frac{\partial g}{\partial q}.
\]
Then the time evolution of any quantity is $\dot f = \{f, H\}$, and position and
momentum satisfy $\{q, p\} = 1$.

What the new form bought was structure, and it paid off over the next ninety years.
Conserved quantities show up on sight: if $H$ does not depend on some $q$, then
$\dot p = -\partial H/\partial q = 0$ and that momentum never changes. Carl Jacobi took
the method up in 1837 and made it the standard tool of celestial mechanics. The flow
in phase space keeps areas fixed, a fact that goes back to Joseph Liouville in 1838,
and Josiah Willard Gibbs built statistical mechanics on it in 1902. In 1926 Erwin
Schrödinger followed Hamilton's own analogy one step further: if the paths of bodies
are to mechanics what rays are to optics, and rays are only an approximation to light
waves, then there should be waves behind the paths of bodies too. And the Poisson
bracket is what Dirac would turn into the quantum commutator.

\begin{question}{What are quaternions, and why do they come up with Hamilton?}
Hamilton is as famous for quaternions as for mechanics. Complex numbers, $a + bi$,
describe rotations in the plane: multiplying by $i$ turns a point by $90^\circ$.
Hamilton spent years looking for numbers that would do the same for rotations in
space. On 16 October 1843, walking with his wife along the Royal Canal in Dublin, he
saw that it needed three imaginary units, not two, and carved the rule into the stone
of Broom Bridge:
\[
  i^2 = j^2 = k^2 = ijk = -1 .
\]
The carving is gone; a plaque marks the place. A quaternion is
$a + bi + cj + dk$, and it does not commute: $ij = k$ but $ji = -k$. It was the first
system of numbers in which the order of a product matters, more than eighty years
before Heisenberg's $xp \neq px$.

They come up in this book for a precise reason. Quaternions of length one describe
every rotation in space, and Chapter~\ref{ch:state} shows that every gate on one
qubit is a rotation of the Bloch sphere. The Pauli matrices, multiplied by $-i$,
obey exactly Hamilton's rule. A single-qubit gate is, up to a phase, a unit
quaternion. That is also why video games and spacecraft use quaternions to store
orientations.
\end{question}

\section{The bridge}

\begin{history}{1925 \textperiodcentered{} Cambridge}
In the autumn of 1925, Paul Dirac, a 23-year-old research student, was reading
Heisenberg's new paper, in which physical quantities no longer commuted: $xp$ was
not equal to $px$. On a Sunday walk he remembered something about Poisson
brackets, but not the details. The libraries were closed, so he waited
impatiently until Monday morning to check. The bracket had the same algebraic
structure as Heisenberg's commutator.
\end{history}

Dirac's rule, now called \emph{canonical quantization}, is a translation table:

\begin{center}
\begin{tabular}{ll}
\hline
classical & quantum \\
\hline
quantities $q$, $p$ & operators $\hat q$, $\hat p$ \\
$\{q, p\} = 1$ & $[\hat q, \hat p] = \hat q\hat p - \hat p\hat q = i\hbar$ \\
$\dot f = \{f, H\}$ & $i\hbar\,\dot{\hat A} = [\hat A, \hat H]$ \\
Hamiltonian $H(q, p)$ & Hamiltonian operator $\hat H(\hat q, \hat p)$ \\
\hline
\end{tabular}
\end{center}

The geometry of phase space does not survive quantization; there are no more
trajectories. The algebra does. The rule is not rigorous for every system, but it
works for every system in this book.

\begin{concept}{The Hamiltonian}
The function, and later the operator, that gives the total energy of a system and
generates its evolution in time.

\textbf{What it fixes:} it puts all the dynamics of a system into a single object,
in a form that survives the move from classical to quantum.

\textbf{What it costs:} you have to know the energy of every interaction in the
system, and for a molecule that is already most of the problem.
\end{concept}

This is why the Hamiltonian is everywhere in quantum algorithms. Hamiltonian
simulation implements $e^{-i\hat H t}$ as a circuit, which is Feynman's original
proposal. Phase estimation extracts the eigenvalues of $\hat H$, which are energies.
Variational algorithms look for the state that minimizes the energy
$\bra{\psi}\hat H\ket{\psi}$. Without the road from Newton to Hamilton, none of
these makes sense.
